% Abhakseite „Das kann ich“ – erzeugt von verzeichnis.py aus den Aufgabendateien
\clearpage
\hypertarget{abhaken}{}
\einheitenkopf*{Das kann ich}
{\small Hake ab, was du kannst.\par}
\medskip
{\small
\par\smallskip\noindent\textbf{Das kennst du schon}\par
\noindent$\square$\quad 1.\ Ich kann Anteile am Streifen ablesen und zeigen.\par
\noindent$\square$\quad 2.\ Ich kann einen Anteil als Bruch schreiben und kürzen.\par
\noindent$\square$\quad 3.\ Ich kann einen Anteil als Bruch schreiben und kürzen – weiter.\par
\noindent$\square$\quad 4.\ Ich kann Brüche und Dezimalzahlen ineinander umwandeln.\par
\noindent$\square$\quad 5.\ Ich kann einen Anteil in Prozent angeben.\par
\noindent$\square$\quad 6.\ Ich kann auf eine Stelle nach dem Komma runden.\par
\noindent$\square$\quad 7.\ Ich kann durch 100 teilen und mit Kommazahlen malnehmen.\par
\noindent$\square$\quad 8.\ Ich kann im Dreisatz erst auf 1 und dann hochrechnen.\par
\noindent$\square$\quad 9.\ Ich kann einen Bruchteil von einer Größe ausrechnen.\par
\noindent$\square$\quad 10.\ Ich finde den Fehler bei einem Anteil.\par
\par\smallskip\noindent\textbf{Einheit 1 von 4 · Prozentsatz berechnen}\par
\noindent$\square$\quad 11.\ Ich erkenne, was das Ganze ist.\par
\noindent$\square$\quad 12.\ Ich kann einen Streifen in gleich große Teile einteilen.\par
\noindent$\square$\quad 13.\ Ich kann am Streifen ablesen, wie viel Prozent ein Teil ist.\par
\noindent$\square$\quad 14.\ Ich kann ausrechnen, wie viel Prozent das sind.\par
\noindent$\square$\quad 15.\ Ich kann ausrechnen, wie viel Prozent das sind – weiter.\par
\noindent$\square$\quad 16.\ Ich kann zu einer Prozentzahl passende Zahlen finden.\par
\noindent$\square$\quad 17.\ Ich kann in einem Text ausrechnen, wie viel Prozent das sind.\par
\noindent$\square$\quad 18.\ Ich finde den Fehler in einer Prozentrechnung.\par
\noindent$\square$\quad 19.\ Ich finde den Fehler in einer Prozentrechnung – weiter.\par
\noindent$\square$\quad 20.\ Ich kann begründen, warum ich durch das Ganze teile.\par
\noindent$\square$\quad 21.\ Ich kann mit Prozenten vergleichen, wer besser trifft.\par
\par\smallskip\noindent\textbf{Einheit 2 von 4 · Prozentwert berechnen}\par
\noindent$\square$\quad 22.\ Ich kann am Streifen ausrechnen, wie viel 50\,\%, 25\,\% oder 10\,\% von einer Größe sind.\par
\noindent$\square$\quad 23.\ Ich kann über 1\,\% ausrechnen, wie viel ein Prozentanteil ist.\par
\noindent$\square$\quad 24.\ Ich kann einen Prozentanteil mit einer Kommazahl ausrechnen.\par
\noindent$\square$\quad 25.\ Ich kann in einem Text ausrechnen, wie viel Euro das sind.\par
\noindent$\square$\quad 26.\ Ich kann in einem Text ausrechnen, wie viel Euro das sind – weiter.\par
\noindent$\square$\quad 27.\ Ich finde den Fehler beim Ausrechnen eines Prozentwerts.\par
\noindent$\square$\quad 28.\ Ich finde den Fehler beim Ausrechnen eines Prozentwerts – weiter.\par
\noindent$\square$\quad 29.\ Ich kann begründen, warum ich zuerst 1\,\% ausrechne.\par
\noindent$\square$\quad 30.\ Ich kann entscheiden, welches Angebot günstiger ist.\par
\par\smallskip\noindent\textbf{Einheit 3 von 4 · Grundwert berechnen}\par
\noindent$\square$\quad 31.\ Ich erkenne, was gegeben und was gesucht ist.\par
\noindent$\square$\quad 32.\ Ich kann am Streifen vom Teil auf das Ganze schließen.\par
\noindent$\square$\quad 33.\ Ich kann aus einem Teil das Ganze ausrechnen.\par
\noindent$\square$\quad 34.\ Ich kann aus einem Teil das Ganze ausrechnen – weiter.\par
\noindent$\square$\quad 35.\ Ich kann unterscheiden, was gesucht ist, und dann richtig rechnen.\par
\noindent$\square$\quad 36.\ Ich finde den Fehler, wenn das Ganze gesucht ist.\par
\noindent$\square$\quad 37.\ Ich kann begründen, warum ich am Ende mal 100 nehme.\par
\noindent$\square$\quad 38.\ Ich kann ausrechnen, wie viel Geld noch fehlt.\par
\par\smallskip\noindent\textbf{Einheit 4 von 4 · Prozentuale Veränderung}\par
\noindent$\square$\quad 39.\ Ich erkenne, ob „um“ oder „auf“ gemeint ist.\par
\noindent$\square$\quad 40.\ Ich erkenne, welcher Wert der alte ist.\par
\noindent$\square$\quad 41.\ Ich kann ausrechnen, wie viel etwas nach einer Erhöhung oder Senkung kostet.\par
\noindent$\square$\quad 42.\ Ich kann ausrechnen, wie viel etwas nach einer Erhöhung oder Senkung kostet – weiter.\par
\noindent$\square$\quad 43.\ Ich kann den neuen Preis in einem Schritt mit einer Kommazahl ausrechnen (kommt in Klasse 9).\par
\noindent$\square$\quad 44.\ Ich kann ausrechnen, um wie viel Prozent sich ein Preis verändert hat.\par
\noindent$\square$\quad 45.\ Ich kann ausrechnen, wie viel etwas vor einer Änderung gekostet hat.\par
\noindent$\square$\quad 46.\ Ich kann Preise mit und ohne Mehrwertsteuer ausrechnen (neu in diesem Jahr).\par
\noindent$\square$\quad 47.\ Ich kann zwei Prozentangaben richtig vergleichen.\par
\noindent$\square$\quad 48.\ Ich kann eine Steigung in Prozent verstehen und ausrechnen (kommt in Klasse 10).\par
\noindent$\square$\quad 49.\ Ich kann ausrechnen, um wie viel Prozent sich ein Preis verändert hat – weiter.\par
\noindent$\square$\quad 50.\ Ich finde den Fehler bei einer Veränderung in Prozent.\par
\noindent$\square$\quad 51.\ Ich kann begründen, welcher Wert 100\,\% ist.\par
\noindent$\square$\quad 52.\ Ich kann entscheiden, ob sich zwei Änderungen aufheben.\par
}
